% Generated draft: requires mathematical and source review.
\documentclass[11pt]{article}
\usepackage[margin=25mm]{geometry}
\usepackage{fontspec}
\setmainfont{Latin Modern Roman}
\usepackage{amsmath,amssymb,mathtools,unicode-math}
\setmathfont{Latin Modern Math}
\usepackage{enumitem,needspace,xurl,hyperref}
\setlength{\parindent}{0pt}
\setlength{\parskip}{5pt}
\setlength{\emergencystretch}{3em}
\newcommand{\N}{\mathbb{N}}
\newcommand{\Z}{\mathbb{Z}}
\newcommand{\Q}{\mathbb{Q}}
\newcommand{\R}{\mathbb{R}}
\newcommand{\C}{\mathbb{C}}
\newcommand{\F}{\mathbb{F}}
\newcommand{\A}{\mathbb{A}}
\newcommand{\PP}{\mathbb{P}}
\newcommand{\E}{\mathbb{E}}
\newcommand{\eps}{\varepsilon}
\newcommand{\dd}{\,\mathrm{d}}
\begin{document}
\section*{UnsolvedMath ID: 30001682}
\textbf{Source problem code:} catalogue entry
\subsection{Definitions and mathematical statement}
Let \(\mathcal T\) be the set of formal variable-free terms generated by the symbol \(1\) and a binary operation \(\star\). Let \(\sim\) be the congruence generated by the left-distributive identity
\[
 A\star(B\star C)\sim(A\star B)\star(A\star C).
\]
The supplied formulation stipulates a family of strictly increasing, nonidentity functions \(f_T:\N\to\N\), indexed by \(T\in\mathcal T\), satisfying
\[
 S\sim T\Longrightarrow f_S=f_T,
 \qquad
 \operatorname{crit}(f_{A\star B})
 =f_A\bigl(\operatorname{crit}(f_B)\bigr),
 \qquad
 \{\operatorname{crit}(f_T):T\in\mathcal T\}=\N.
\]
Under the usual interpretation, the critical point of such a function is
\[
 \operatorname{crit}(f)=\min\{k\in\N:f(k)\ne k\}.
\]
This interpretation is not explicitly defined in the supplied record and requires confirmation. The recursion below evaluates functions at \(0\), so it requires \(0\in\N\), or an explicitly supplied extension of their domains; the record does not state its convention for \(\N\).

Define, on formal terms,
\[
 \phi(1)=0,
 \qquad
 \phi(A\star B)=\max\{\phi(A),f_A(\phi(B))\}.
\]
For positive integers \(m\), write the right-associated powers as
\[
 R_1=1,\qquad R_{m+1}=1\star R_m.
\]
The original statement describes \(A_{n,p}\) using \emph{iterated right divisors} of \(R_{p+1}\); the cleaned statement says only \emph{right divisors}. Retaining the original qualification, its intended description is
\[
 A_{n,p}
 =\{T\in\operatorname{RDiv}_{\mathrm{it}}(R_{p+1}):\phi(T)\le n\}.
\]
Here \(\operatorname{RDiv}_{\mathrm{it}}\) is a placeholder for the source's undefined iterated-right-divisor relation, not a completed mathematical definition. The record does not specify whether divisibility is taken among formal terms, modulo \(\sim\), or in another algebraic object. If equivalence classes are intended, the use of \(\phi\) on those classes also needs justification.

The proposed assertion, for the intended function family and intended divisor and F\o lner definitions, is
\[
 \forall\varepsilon\in\R_{>0}\ \exists p_0\in\N\
 \forall p\in\N\ \Bigl(p>p_0\Longrightarrow
 \exists n_0\in\N\ \forall n\in\N\
 \bigl(n\ge n_0\Longrightarrow
 A_{n,p}\text{ is }\varepsilon\text{-F\o lner}\bigr)\Bigr).
\]
Thus the finite exceptional set of indices \(n\) may depend on both \(\varepsilon\) and \(p\). The supplied background calls these sets finite. However, the record supplies no ambient group, semigroup, or action, no transformations against which the F\o lner condition is tested, and no boundary inequality or normalization defining \(\varepsilon\)-F\o lner. Consequently this display preserves the proposed quantifiers but remains an incomplete formulation pending those definitions. The record also does not specify whether the assertion concerns one distinguished function family or every family satisfying the displayed conditions.
\subsection{Short English statement}
The question asks whether sufficiently large right-associated powers of the generator yield, through bounded-\(\phi\) iterated right divisors, \(\varepsilon\)-F\o lner sets for every sufficiently large bound \(n\). The threshold for \(n\) may depend on \(\varepsilon\) and the chosen power. The intended divisor relation, ambient action, precise F\o lner test, and function-family quantification are missing, so the supplied record does not determine a complete precise conjecture.
\subsection{Sources}
\begin{itemize}
\item Local UnsolvedMath registry, \texttt{data/unsolvedmath/problems.json}, permanent ID \(30001682\), title ``F\o lner Sets from Variable-Free Terms,'' exact source code \texttt{catalogue entry}. The supplied record attributes the problem to the Oberwolfach Reports Open Problems collection and gives proposed year \(2011\); no individual proposer is supplied.
\item \emph{Set Theory} (2011), Oberwolfach Reports, pp.\ 85--140, DOI \texttt{10.4171/owr/2011/02}. Supplied report citation and literature-source label ``Oberwolfach problem statement'': \url{https://doi.org/10.4171/owr/2011/02}.
\item Supplied literature-source label ``current amenability survey'': \url{https://arxiv.org/abs/2305.07113}. The record supplies no author, title, or passage identifying its relevance to the exact conjecture.
\item Literature assessment date: \texttt{2026-08-21}. The supplied dated curation labels the problem open and reports that no later proof of the exact conjectured test was found. This is supplied dated curation, not an independently verified literature assessment.
\item The registry labels the statement \texttt{corrected\_verified} and reports that the official report was checked and the binary operation, function system, critical-point identities, and F\o lner-set definition were restored. This is supplied verification metadata, not independent verification here; the supplied statement still omits an explicit F\o lner-set definition.
\end{itemize}
\end{document}
